% ---------------------------------------------------------------------------
% model_notes.tex
% Appendix: supplementary modelling notes.
% This file is \input by main.tex after \appendix; it is not a standalone
% document.
% ---------------------------------------------------------------------------

\section{Supplementary Model Notes}
\label{app:notes}

\subsection{Why the prediction is conditional on the quarterfinal stage}
\label{app:conditional}

A pre-tournament forecast must average over an enormous outcome space: group
compositions, group-stage results, and several knockout rounds, each of which
reshuffles who meets whom. Once the tournament reaches the quarterfinals,
almost all of that uncertainty has been \emph{resolved}. The sample space
collapses to the eight surviving teams and the seven remaining matches, and
the realized bracket---including which half each team occupies---is known
exactly. Conditioning on this information has three consequences.

First, eliminated teams are excluded by construction: a team knocked out in
the group stage or an earlier round has conditional championship probability
zero, and the model never represents it. Second, the conditional
probabilities of the eight quarterfinalists necessarily sum to one, which
makes the output directly interpretable as a complete probability
distribution over the champion. Third, the conditional forecast is generally \emph{sharper} than a
pre-tournament one: for a team that was expected to struggle, having already
reached the last eight is strong evidence of current quality, which enters the
model through the updated covariates (Elo after the earlier rounds, tournament
form, rest days). Conversely, quantities from before the tournament---title
odds, group-stage simulations---answer a different question and must not be
compared numerically to the conditional probabilities computed here.

\subsection{Why Monte Carlo simulation is useful with only seven matches
remaining}
\label{app:whymc}

With eight teams and seven binary match outcomes, the remaining tournament has
only $2^7 = 128$ possible trajectories, and the championship probabilities
admit the closed-form factorization of Appendix~\ref{app:path}. Monte Carlo
simulation is nevertheless the method of choice, for four reasons.

\begin{enumerate}
  \item \textbf{Extensibility.} The closed form exists only because match
        outcomes are conditionally independent given fixed parameters. Every
        realistic enrichment breaks it: drawing model parameters from a
        Bayesian posterior in each replication, injecting random injury or
        fatigue shocks, or letting a simulated extra-time semifinal reduce a
        finalist's rest covariate all induce dependence across matches.
        The simulator of Algorithm~\ref{alg:tournament} absorbs each of these
        changes with a few added lines; the analytic formula has to be
        re-derived or abandoned.
  \item \textbf{Richer outputs.} The same simulated trajectories that
        estimate champion probabilities also deliver, at no extra cost, the
        distribution of final pairings, the probability that a given team
        reaches each round, expected goals per fixture, and the joint law of
        any bracket event of interest.
  \item \textbf{Negligible cost and controlled error.} With $B = 100{,}000$
        replications the worst-case standard error is $0.16$ percentage
        points, far below model uncertainty, and the run completes in about a
        second. Increasing $B$ buys arbitrary accuracy at linear cost.
  \item \textbf{Transparency.} The code mirrors the tournament's structure
        line by line (four quarterfinals, two semifinals, one final), which
        makes it easy to audit against the bracket---often harder for a
        hand-derived sum over 128 paths.
\end{enumerate}

\subsection{How the bracket path affects championship probability}
\label{app:path}

Championship probability is not a function of team strength alone; it depends
on \emph{who can be met, and when}. Consider a team $t$ with quarterfinal
opponent $u$, let $\mathcal{V}$ denote the two teams in the paired
quarterfinal (the potential semifinal opponents), and let $\mathcal{W}$ denote
the four teams in the other half of the bracket (the potential final
opponents). Because match outcomes are conditionally independent given the
parameters, the conditional championship probability factorizes along the
path:
\begin{equation}
\Prob(C = t \mid \mathcal{Q})
\;=\;
\underbrace{p_{tu}}_{\text{quarterfinal}}
\;\times\;
\underbrace{\sum_{v \in \mathcal{V}} q_v\, p_{tv}}_{\text{semifinal}}
\;\times\;
\underbrace{\sum_{w \in \mathcal{W}} f_w\, p_{tw}}_{\text{final}},
\label{eq:pathfactorization}
\end{equation}
where $q_v$ is the probability that $v$ wins the paired quarterfinal, and
$f_w$ is the probability that $w$ emerges from the other half, computed
recursively as $f_w = \Prob(w \text{ wins its QF}) \sum_{u'} q_{u'}\,p_{wu'}$
with the sum running over the other half's adjacent quarterfinal (so that
$\sum_{w \in \mathcal{W}} f_w = 1$). Equation
\eqref{eq:pathfactorization} is exactly what the Monte Carlo procedure
estimates, and it makes the bracket dependence explicit: the semifinal factor
for France averages over \{Spain, Belgium\}, whereas the semifinal factor for
England averages over \{Argentina, Switzerland\}.

Under the estimated inputs this matters visibly. France enters the
quarterfinals with a higher Elo rating than England (2207 versus 2150), yet
England's championship probability (0.145) exceeds France's (0.133). The
reason is France's path: it must first break down Morocco, the best-fitted
defence among the eight, and then---because the top half's two strongest
teams, France and Spain, sit in \emph{adjacent} quarterfinals---most likely
face Spain already in semifinal~1, so at most one of the two can reach the
final. A counterfactual bracket that placed Spain in the bottom half would
raise France's title probability without changing a single team parameter.
This is why re-seeding debates are substantive: the draw itself carries
probability mass.

\subsection{Extensions: expected goals, injuries, betting odds, and Bayesian
uncertainty}
\label{app:extensions}

\paragraph{Expected goals (xG).}
Raw goal counts are noisy realizations of chance creation. A first upgrade
fits the attack and defence parameters of \eqref{eq:lograte} to xG totals
rather than goals, reducing finishing noise in estimation. A second adds an
xG-differential covariate to the log-rate,
$\log \lambda_{ij} \mathrel{+}= \kappa\,(X_i - X_j)$, where $X_i$ is the
team's tournament xG difference per match; this captures underlying
performance that results-based measures like Elo absorb only slowly.

\paragraph{Injuries and suspensions.}
Player availability can be folded in through an availability-adjusted rating
$E_i^{\mathrm{adj}} = E_i - \pi_i$, where $\pi_i$ converts missing players
into rating points (e.g.\ minutes-weighted contribution estimates). More
ambitiously, availability can be made stochastic: in each Monte Carlo
replication, draw an injury/suspension scenario for the remaining rounds and
adjust the affected team's parameters before simulating---one of the
dependence-inducing extensions for which the closed form of
Appendix~\ref{app:path} no longer applies.

\paragraph{Betting odds.}
Bookmaker prices aggregate large information sets. After removing the
bookmaker margin (overround), match odds imply advancement probabilities
$p_{ij}^{\mathrm{mkt}}$ that can be used in three ways: as a benchmark to
validate the model's $p_{ij}$; as a blend
$\tilde{p}_{ij} = \omega\, p_{ij}^{\mathrm{mkt}} + (1 - \omega)\, p_{ij}$
with $\omega$ chosen by historical calibration; or as calibration targets to
which the coefficients $(\gamma, \delta, \rho, \eta)$ are tuned.

\paragraph{Bayesian uncertainty.}
The pipeline treats $\theta = (\alpha, \{a_i\}, \{d_i\}, \gamma, \delta,
\rho, \eta)$ as known, so Table~\ref{tab:output} understates real
uncertainty. A Bayesian treatment places priors on $\theta$, obtains the
posterior $p(\theta \mid \mathcal{D})$ from historical match data
$\mathcal{D}$ (e.g.\ via MCMC in Stan or PyMC), and reports the posterior
predictive championship probability
\begin{equation}
\Prob(C = t \mid \mathcal{Q}, \mathcal{D})
\;=\;
\int \Prob(C = t \mid \mathcal{Q}, \theta)\,
p(\theta \mid \mathcal{D})\, d\theta,
\label{eq:posteriorpredictive}
\end{equation}
implemented by drawing a fresh $\theta^{(b)}$ from the posterior in each
Monte Carlo replication before simulating the bracket. The spread of
$\Prob(C = t \mid \mathcal{Q}, \theta^{(b)})$ across draws yields credible
intervals for each championship probability---typically wider, and more
honest, than the Monte Carlo error bars alone. Time-decayed likelihood
weights in the spirit of Dixon--Coles integrate naturally into the same
framework.

\section{Data Provenance and Reproducibility}
\label{app:repro}

This appendix records everything needed to audit or reproduce the numbers in
this note. All inputs derive from one public match-level dataset; one
published rating table is used solely for cross-validation.

\subsection{Sources and retrieval}
\label{app:sources}

\begin{itemize}
  \item \textbf{Match results.} \texttt{results.csv} from the
        community-maintained repository of \cite{martj42results}
        (\url{https://github.com/martj42/international_results}, raw file
        \texttt{master/results.csv}), retrieved on 9~July 2026:
        49{,}505 international matches from 30~November 1872 through
        7~July 2026, plus the four scheduled quarterfinal fixtures. The file
        contains all 96 matches of the 2026 World Cup played through the
        round of~16 (72 group matches, 16 in the round of~32, 8 in the round
        of~16). MD5 checksum:
        \texttt{3820fec01802a4303d91b00d6e298bed}.
  \item \textbf{Penalty shoot-outs.} \texttt{shootouts.csv} from the same
        repository and retrieval date: 681 shoot-outs, 1967--2026. MD5
        checksum: \texttt{3461b3d538309f7f17b63259e4b0b524}.
  \item \textbf{Cross-validation only.} The published Elo table of 7~July
        2026 \cite{eloratings} was used to verify the Elo replication of
        Section~\ref{sec:data} (ranking and pairwise gaps of the eight
        quarterfinalists); it enters no computation.
\end{itemize}
Both raw files are archived with the project sources, so every number in
this note can be reproduced offline; the checksums allow byte-level
verification against any fresh download.

\subsection{Processing rules}
\label{app:processing}

\begin{enumerate}
  \item Rows with missing scores (unplayed fixtures) are excluded from all
        estimation; matches are processed in chronological order.
  \item Recorded scores include extra-time goals and exclude penalty
        shoot-outs (dataset convention). Matches decided by shoot-out
        therefore enter both the Elo replay and the regression as draws.
  \item \emph{Elo replay:} every team starts at 1500; after each match,
        $R \leftarrow R + K G (W - W_e)$ with expected score
        $W_e = (1 + 10^{-\mathrm{dr}/400})^{-1}$, where $\mathrm{dr}$ is the
        rating difference plus 100 points for the home side at non-neutral
        venues; $K = 60$ (World Cup finals), $50$ (continental finals),
        $40$ (qualifiers, Nations League), $30$ (other tournaments), $20$
        (friendlies); goal-difference multiplier $G = 1$, $1.5$, or
        $(11 + |\Delta|)/8$ for winning margins $\le 1$, $= 2$, $\ge 3$.
  \item \emph{Poisson regression:} all matches from 1~January 2016 through
        7~July 2026 (10{,}048 matches, two observations each, 294 teams);
        time-decay weights $0.5^{\Delta t / 1095\,\mathrm{days}}$ with
        $\Delta t$ measured to 7~July 2026 (three-year half-life; effective
        sample size 7{,}891); ridge penalty of $1.0$ on the attack/defence
        dummies only; attack/defence recentred to mean zero after
        convergence; Elo differences scaled by $1/100$ inside the solver
        and rescaled on output; a home-advantage dummy is included in
        estimation (fitted $0.2342$) and switched off in the neutral-venue
        knockout predictions.
  \item \emph{Covariates:} form $F_i$ is the mean goal difference over the
        team's previous five internationals; rest $R_i$ is the number of
        days since the team's previous international, capped at 14.
        Identical definitions are used in estimation and prediction.
  \item \emph{Tie-break:} $\eta$ is fitted by Newton--Raphson maximum
        likelihood on the 681 shoot-outs, using each side's pre-match
        replicated Elo rating.
\end{enumerate}

\subsection{Reproduction steps}
\label{app:steps}

With Python $\ge 3.9$, \texttt{numpy}, and \texttt{pandas}:

\begin{verbatim}
# (optional) refresh the data; otherwise the archived copies are used
curl -sL -O https://raw.githubusercontent.com/martj42/\
international_results/master/results.csv
curl -sL -O https://raw.githubusercontent.com/martj42/\
international_results/master/shootouts.csv

python3 fit_parameters.py      # Elo replay + IRLS + eta fit  (~2 s)
python3 simulate_worldcup.py   # 100,000 simulations, seed 42 (~1 s)
\end{verbatim}

\noindent
\texttt{fit\_parameters.py} regenerates \texttt{data\_template.csv}
(Table~\ref{tab:inputs}) and prints the fitted coefficients
\[
\alpha = 0.150326,\qquad
\gamma = 0.000972,\qquad
\delta = 0.005260,\qquad
\rho = -0.000886,\qquad
\eta = 0.001382,
\]
which are hard-coded in \texttt{simulate\_worldcup.py}. Every stage of the
pipeline is deterministic---the simulation uses
\texttt{numpy.random.default\_rng(42)}---so independent runs on the archived
data reproduce Table~\ref{tab:output} exactly.
